% ============================================================
% 03_arquitectura.tex
% Responsable: Integrante X
% ============================================================

\section{Análisis de Arquitectura}

\subsection{Visión general}

% TODO: Insertar diagrama de arquitectura C-Lock (figura del paper o redibujada)
% \begin{figure}[H]
%     \centering
%     \includegraphics[width=0.8\textwidth]{figures/arquitectura_clock.pdf}
%     \caption{Arquitectura del sistema C-Lock.}
%     \label{fig:arch}
% \end{figure}

\subsection{Flujo de operación}

El flujo de ejecución cuando un core intenta adquirir un lock es el siguiente:

\begin{enumerate}
    \item El core ejecuta \texttt{BEGIN\_C\_LOCK(lock\_id)}.
    \item El \textit{Conflict Checker} del \textit{Pool\_i} correspondiente consulta
          al C-LockManager si el lock está ocupado.
    \item Si existe conflicto: el clock-gating se activa y el core se suspende.
    \item Cuando el lock se libera (\texttt{END\_C\_LOCK}), el C-LockManager notifica
          al core en espera y reactiva su reloj.
    \item El core retoma la ejecución y entra a la sección crítica.
\end{enumerate}

\subsection{Decisiones de diseño relevantes}

% TODO: Analizar trade-offs del diseño:
%   - ¿Por qué hardware y no software?
%   - ¿Por qué clock-gating y no power-gating?
%   - Granularidad de locks (por lock_id vs global)
%   - Overhead de área hardware (costo en transistores/área)

\subsection{Comparación con otros mecanismos}

\begin{table}[H]
    \centering
    \caption{Comparación de mecanismos de sincronización.}
    \begin{tabular}{lccc}
        \toprule
        Mecanismo       & Consumo energético & Overhead SW & Complejidad HW \\
        \midrule
        Spin-lock       & Alto               & Bajo        & Ninguna        \\
        Mutex (OS)      & Medio              & Alto        & Ninguna        \\
        C-Lock          & Bajo               & Mínimo      & Moderada       \\
        \bottomrule
    \end{tabular}
    \label{tab:comparacion}
\end{table}

% TODO: Completar tabla con datos cuantitativos del paper.
